% =====================================================================
% บทที่ 4 การวิเคราะห์ระบบ และพัฒนาโปรแกรม (คู่มือข้อ 17.4)
%   4.1 การวิเคราะห์ระบบ   4.2 การออกแบบระบบ
%   4.3 การพัฒนาโปรแกรม   4.4 การทดสอบระบบ
% หัวข้อย่อยปรับได้ตามลักษณะของโครงงาน (งานวิจัย/ฮาร์ดแวร์ ดูคู่มือข้อ 17.4)
% *** ก่อนส่ง: ลบข้อความตัวอย่าง (...) และภาพตัวอย่าง (กรอบสีเทา) ***
% =====================================================================

\section{การวิเคราะห์ระบบ}
\label{sec:analysis}

\subsection{ผังแสดงทิศทางของกิจกรรม}

การทำงานของระบบเดิมและเอกสารที่เกี่ยวข้องแสดงดังภาพที่~\ref{fig:activity}

% ภาพตัวอย่าง (placeholder): แทนด้วย \includegraphics[width=0.8\linewidth]{images/activity.png}
\begin{figure}[H]
\centering
\cskkuPlaceholder[ผังแสดงทิศทางของกิจกรรม (Activity diagram)]{0.8\linewidth}{5cm}
\caption{ผังแสดงทิศทางของกิจกรรมของระบบเดิม}
\label{fig:activity}
\end{figure}

\subsection{รายการความต้องการ}

ความต้องการเชิงฟังก์ชันของระบบแสดงดังตารางที่~\ref{tab:fr}

\begin{table}[H]
\caption{ความต้องการเชิงฟังก์ชันของระบบ}
\label{tab:fr}
\centering
\begin{tabular}{|C{1.6cm}|L{4cm}|L{\dimexpr\linewidth-5.6cm-6\tabcolsep-4\arrayrulewidth\relax}|}
\hline
\textbf{รหัส} & \textbf{ความต้องการ} & \textbf{รายละเอียด} \\ \hline
FR-01 & สมัครสมาชิก & ผู้ใช้สมัครสมาชิกด้วยอีเมลของมหาวิทยาลัย \\ \hline
FR-02 & เข้าสู่ระบบ & ผู้ใช้เข้าสู่ระบบด้วยอีเมลและรหัสผ่าน \\ \hline
FR-03 & \ldots & ระบบต้องสามารถ\ldots{} \\ \hline
\end{tabular}
\end{table}

ความต้องการที่ไม่ใช่เชิงฟังก์ชัน มีดังนี้
\begin{enumerate}
  \item ด้านประสิทธิภาพ: ระบบแสดงผลแต่ละหน้าภายใน\ldots{} วินาที
  \item ด้านความปลอดภัย: ระบบเก็บรหัสผ่านด้วยการเข้ารหัสแบบ\ldots{}
  \item ด้านการใช้งาน: \ldots{}
\end{enumerate}

\subsection{Use case diagram}

ผู้ใช้และกรณีการใช้งานของระบบแสดงดังภาพที่~\ref{fig:usecase}
และคำอธิบายกรณีการใช้งาน (Use case text) แสดงดังตารางที่~\ref{tab:uc01}

\begin{figure}[H]
\centering
\cskkuPlaceholder[Use case diagram]{0.8\linewidth}{6cm}
\caption{Use case diagram ของระบบ}
\label{fig:usecase}
\end{figure}

\begin{table}[H]
\caption{คำอธิบายกรณีการใช้งาน UC-01 เข้าสู่ระบบ}
\label{tab:uc01}
\centering
\begin{tabular}{|L{3.6cm}|L{\dimexpr\linewidth-3.6cm-4\tabcolsep-3\arrayrulewidth\relax}|}
\hline
\textbf{Use case} & UC-01 เข้าสู่ระบบ \\ \hline
\textbf{ผู้กระทำ (Actor)} & นักศึกษา อาจารย์ที่ปรึกษา \\ \hline
\textbf{เงื่อนไขก่อนเริ่ม} & ผู้ใช้สมัครสมาชิกแล้ว \\ \hline
\textbf{ขั้นตอนหลัก} & 1) ผู้ใช้กรอกอีเมลและรหัสผ่าน \newline
                       2) ระบบตรวจสอบข้อมูล \newline
                       3) ระบบแสดงหน้าหลักตามบทบาทของผู้ใช้ \\ \hline
\textbf{ขั้นตอนทางเลือก} & 2a) ข้อมูลไม่ถูกต้อง ระบบแสดงข้อความแจ้งเตือน \\ \hline
\textbf{ผลลัพธ์} & ผู้ใช้เข้าสู่ระบบสำเร็จ \\ \hline
\end{tabular}
\end{table}

\subsection{System sequence diagram}

ภาพรวมการทำงานของระบบแสดงด้วย System sequence diagram ดังภาพที่\ldots{}

\subsection{Conceptual class diagram}

ความสัมพันธ์เชิงแนวคิดของข้อมูลในระบบแสดงดังภาพที่\ldots{}

\section{การออกแบบระบบ}
\label{sec:design}

\subsection{สถาปัตยกรรมของระบบ}

ระบบแบ่งเป็น 3 ส่วน ได้แก่ ส่วนหน้า (Frontend) ส่วนหลัง (Backend) และฐานข้อมูล
ดังภาพที่~\ref{fig:architecture} ส่วนหน้าติดต่อกับส่วนหลังผ่าน REST API~\cite{fielding2000}

\begin{figure}[H]
\centering
\begin{tikzpicture}[
    node distance=10bp and 34bp,
    tier/.style={draw=black!70, fill=black!4, rounded corners=3pt, align=center,
      font=\normalfont\small, minimum height=46bp, text width=78bp},
    arr/.style={{Stealth[length=5bp]}-{Stealth[length=5bp]}, line width=0.8pt, draw=black!60},
    lbl/.style={font=\normalfont\footnotesize, inner sep=2bp}]
  \node[tier] (fe) {ส่วนหน้า\\(React)};
  \node[tier, right=of fe] (be) {ส่วนหลัง\\(Node.js + Express)};
  \node[tier, right=of be] (db) {ฐานข้อมูล\\(PostgreSQL)};
  \draw[arr] (fe) -- node[lbl, above]{REST API} (be);
  \draw[arr] (be) -- node[lbl, above]{SQL} (db);
\end{tikzpicture}
\caption{สถาปัตยกรรมของระบบ}
\label{fig:architecture}
\end{figure}

\subsection{Class diagram}

รายละเอียดคลาสและความสัมพันธ์แสดงดังภาพที่\ldots{}

\subsection{Sequence diagram}

พฤติกรรมเชิงพลวัตของระบบแสดงด้วย Sequence diagram ดังภาพที่\ldots{}

\subsection{การออกแบบฐานข้อมูล}

ฐานข้อมูลประกอบด้วย\ldots{} ตาราง ความสัมพันธ์แสดงด้วย E-R diagram ดังภาพที่\ldots{}
โครงสร้างของตาราง users แสดงดังตารางที่~\ref{tab:db-users}

% ตารางยาวหลายหน้า: ใช้ longtable ชื่อตารางจะพิมพ์ซ้ำพร้อม "(ต่อ)" ให้อัตโนมัติ
\begin{longtable}{|L{3.4cm}|L{3cm}|C{1.6cm}|L{\dimexpr\linewidth-8cm-8\tabcolsep-5\arrayrulewidth\relax}|}
\caption{โครงสร้างตาราง users}\label{tab:db-users}\\
\hline
\textbf{ชื่อฟิลด์} & \textbf{ชนิดข้อมูล} & \textbf{คีย์} & \textbf{คำอธิบาย} \\ \hline
\endfirsthead
\caption[]{โครงสร้างตาราง users~\cskkuContinued}\\
\hline
\textbf{ชื่อฟิลด์} & \textbf{ชนิดข้อมูล} & \textbf{คีย์} & \textbf{คำอธิบาย} \\ \hline
\endhead
\cskkuCode{id}            & INTEGER      & PK & รหัสผู้ใช้ \\ \hline
\cskkuCode{email}         & VARCHAR(100) & UQ & อีเมลของมหาวิทยาลัย \\ \hline
\cskkuCode{password_hash} & VARCHAR(255) &    & รหัสผ่านที่เข้ารหัสแล้ว \\ \hline
\cskkuCode{role}          & VARCHAR(20)  &    & บทบาทของผู้ใช้ \\ \hline
\cskkuCode{created_at}    & TIMESTAMP    &    & วันและเวลาที่สร้างบัญชี \\ \hline
\end{longtable}

\subsection{ต้นแบบส่วนติดต่อกับผู้ใช้}

ต้นแบบส่วนติดต่อกับผู้ใช้ (Mockup user interface) ของหน้าหลักแสดงดังภาพที่\ldots{}
และภาพหน้าจอของระบบทั้งหมดแสดงในภาคผนวก~ก

\section{การพัฒนาโปรแกรม}
\label{sec:implementation}

\subsection{ภาษาและเครื่องมือที่ใช้ในการพัฒนา}

ระบบพัฒนาด้วยเครื่องมือในตารางที่~\ref{tab:software} โดยใช้ไลบรารี\ldots{}

\subsection{โครงสร้างของโปรแกรม}

โปรแกรมส่วนหลังแบ่งเป็นโมดูลตามหน้าที่ ตัวอย่างฟังก์ชันตรวจสอบการเข้าสู่ระบบ
แสดงดังโปรแกรมที่~\ref{lst:login}

% โค้ดโปรแกรม: lstlisting (พิมพ์ภาษาไทยในความเห็นได้)
% หรือนำเข้าจากไฟล์: \lstinputlisting[language=Java, caption={...}, label={...}]{code/Main.java}
\begin{lstlisting}[language=Java, caption={ฟังก์ชันตรวจสอบการเข้าสู่ระบบ}, label={lst:login}]
public User login(String email, String pw) {
    // ค้นหาผู้ใช้จากอีเมล (ไม่สนตัวพิมพ์เล็ก-ใหญ่)
    User u = users.findByEmail(email.toLowerCase());
    if (u == null || !hasher.verify(pw, u)) {
        throw new AuthException(
            "อีเมลหรือรหัสผ่านไม่ถูกต้อง");
    }
    return u;
}
\end{lstlisting}

\subsection{ข้อกำหนดเฉพาะของแต่ละส่วน}

ข้อกำหนดเฉพาะของโปรแกรม (Program specification) แต่ละส่วนมีดังนี้

% หัวข้อย่อยระดับลึก (คู่มือหน้า 24): \subsubsection -> (1)  \paragraph -> (1.1)
% (ตัวเลือก plainnumbering: 1.1.1 และ (1)) ไม่ควรลึกกว่านี้ ถ้าต้องแบ่งย่อยอีกให้ใช้ enumerate
% สารบัญแสดงถึงหัวข้อระดับ 4.3.3 เท่านั้น ตามคู่มือข้อ 13
\subsubsection{ส่วนสมัครสมาชิก}

รับข้อมูล\ldots{} ตรวจสอบ\ldots{} แล้วบันทึกลงตาราง\ldots{}

\subsubsection{ส่วนจัดการข้อมูล\ldots{}}

\paragraph{การเพิ่มข้อมูล}

รับข้อมูล\ldots{} จากแบบฟอร์ม\ldots{}

\paragraph{การแก้ไขข้อมูล}

ค้นหาข้อมูลจาก\ldots{} แล้วแก้ไข\ldots{}

\subsection{การติดต่อกับฐานข้อมูลและบริการภายนอก}

ระบบติดต่อกับฐานข้อมูลผ่าน\ldots{}

\section{การทดสอบระบบ}
\label{sec:testing}

\subsection{วิธีการทดสอบ}

ทดสอบการยอมรับของผู้ใช้ (User acceptance testing) กับผู้ใช้จำนวน\ldots{} คน
ตามกรณีทดสอบในตารางที่~\ref{tab:uat} รายละเอียดของแต่ละกรณีทดสอบเขียนตามรูปแบบตารางที่~\ref{tab:tc01}

\begin{table}[H]
\caption{รายละเอียดกรณีทดสอบ TC-01}
\label{tab:tc01}
\centering
\begin{tabular}{|L{3.6cm}|L{\dimexpr\linewidth-3.6cm-4\tabcolsep-3\arrayrulewidth\relax}|}
\hline
\textbf{รหัสกรณีทดสอบ} & TC-01 \\ \hline
\textbf{ชื่อกรณีทดสอบ} & เข้าสู่ระบบด้วยข้อมูลที่ถูกต้อง \\ \hline
\textbf{เงื่อนไขก่อนทดสอบ} & มีบัญชีผู้ใช้ \cskkuCode{test@kku.ac.th} ในระบบ \\ \hline
\textbf{ขั้นตอน} & 1) เปิดหน้าเข้าสู่ระบบ \newline
                   2) กรอกอีเมล \cskkuCode{test@kku.ac.th} และรหัสผ่าน \cskkuCode{P@ssw0rd} \newline
                   3) คลิกปุ่มเข้าสู่ระบบ \\ \hline
\textbf{ผลที่คาดหวัง} & ระบบแสดงหน้าหลักพร้อมชื่อผู้ใช้ \\ \hline
\textbf{ผลการทดสอบ} & ผ่าน \\ \hline
\end{tabular}
\end{table}

\subsection{ผลการทดสอบ}

ผลการทดสอบตามกรณีทดสอบแสดงดังตารางที่~\ref{tab:uat}

\begin{table}[H]
\caption{ผลการทดสอบการยอมรับของผู้ใช้}
\label{tab:uat}
\centering
\begin{tabular}{|C{1.6cm}|L{5.2cm}|L{\dimexpr\linewidth-8.6cm-8\tabcolsep-5\arrayrulewidth\relax}|C{1.8cm}|}
\hline
\textbf{รหัส} & \textbf{กรณีทดสอบ} & \textbf{ผลที่คาดหวัง} & \textbf{ผล} \\ \hline
TC-01 & เข้าสู่ระบบด้วยข้อมูลที่ถูกต้อง & แสดงหน้าหลัก & ผ่าน \\ \hline
TC-02 & เข้าสู่ระบบด้วยรหัสผ่านผิด & แสดงข้อความแจ้งเตือน & ผ่าน \\ \hline
TC-03 & \ldots & \ldots & ผ่าน \\ \hline
\end{tabular}
\end{table}

ผลการประเมินความพึงพอใจของผู้ใช้จำนวน\ldots{} คน แสดงดังตารางที่~\ref{tab:satisfaction}

\begin{table}[H]
\caption{ผลการประเมินความพึงพอใจของผู้ใช้}
\label{tab:satisfaction}
\centering
\begin{tabular}{|L{\dimexpr\linewidth-6.6cm-8\tabcolsep-5\arrayrulewidth\relax}|C{1.8cm}|C{1.8cm}|C{3cm}|}
\hline
\textbf{หัวข้อการประเมิน} & $\boldsymbol{\bar{x}}$ & \textbf{S.D.} & \textbf{ระดับความพึงพอใจ} \\ \hline
ความง่ายในการใช้งาน & 4.52 & 0.51 & มากที่สุด \\ \hline
ความถูกต้องของข้อมูล & 4.31 & 0.60 & มาก \\ \hline
ความเร็วในการทำงาน & \ldots & \ldots & \ldots \\ \hline
\textbf{รวม} & \textbf{\ldots} & \textbf{\ldots} & \textbf{\ldots} \\ \hline
\end{tabular}
\end{table}
